\lysh{12683}{4}

\begin{alist}
\item Γνωρίζουμε ότι ο όγκος της δεξαμενής ισούται με το γινόμενο των τριών διαστάσεών της. Επειδή η δεξαμενή έχει βάση τετράγωνη θέτουμε $x$ το μήκος και το πλάτος της οπότε το ύψος της θα είναι $\dfrac{x}{4}$. Άρα ο όγκος της δεξαμενής $V$ θα είναι:

\[V = x \cdot x \cdot \dfrac{x}{4} = \dfrac{x^3}{4}\]

Αφού η δεξαμενή έχει όγκο $V = 16 m^3$ θα έχουμε:

\[V = 16 \Leftrightarrow \dfrac{x^3}{4} = 16 \Leftrightarrow x^3 = 64 \Leftrightarrow x = \sqrt[3]{64} \Leftrightarrow x = \sqrt[3]{4^3} \Leftrightarrow x = 4m.\]

Οπότε η δεξαμενή έχει μήκος και πλάτος ίσα με $4m$ και ύψος ίσο με $1m$.

\item Έστω $x, (x > 0)$ το μήκος της δεξαμενής. Τότε το πλάτος της θα είναι $x-2$ και ο όγκος της δεξαμενής θα ισούται με $V = 2x(x - 2)$.
Οπότε έχουμε:
\begin{align*}
V = 16 \ &\Leftrightarrow \\
2x(x - 2) = 16 \ &\Leftrightarrow \\
x(x - 2) = 8 \ &\Leftrightarrow \\
x^2 - 2x = 8 \ &\Leftrightarrow \\
x^2 - 2x - 8 &= 0
\end{align*}

Έχουμε $\Delta = \beta^2 - 4 \cdot \alpha \cdot \gamma = (-2)^2 - 4 \cdot 1 \cdot (-8) = 4 + 32 = 36 > 0$.
Τότε:
\[x_{1,2} = \dfrac{-\beta \pm \sqrt{\Delta}}{2 \cdot \alpha} \Leftrightarrow x_{1,2} = \dfrac{2 \pm \sqrt{36}}{2 \cdot 1} \Leftrightarrow x_{1,2} = \dfrac{2 \pm 6}{2}\]

\[x_{1,2} = \begin{cases} \dfrac{2 - 6}{2} \\ \dfrac{2 + 6}{2} \end{cases} \Leftrightarrow x_{1,2} = \begin{cases} -2, & \text{απορρίπτεται, γιατί πρέπει } x > 0 \\ 4, & \text{δεκτή} \end{cases}\]

Άρα το μήκος της δεξαμενής είναι $4m$ και το πλάτος $2m$.

\item Αφού η νέα δεξαμενή περιέχει $10m^3$ πετρέλαιο και η βάση της έχει μήκος $4m$ και πλάτος $2m$, αν $x$ είναι το ύψος του υγρού μέσα στη δεξαμενή ο όγκος του υγρού θα είναι:

\[V_{\pi\varepsilon\tau\rho.} = 10 \Leftrightarrow 4 \cdot 2 \cdot x = 10 \Leftrightarrow 8 \cdot x = 10 \Leftrightarrow x = \dfrac{10}{8} \Leftrightarrow x = \dfrac{5}{4}m.\]

Άρα το ύψος του υγρού στη δεξαμενή είναι $\dfrac{5}{4}m$.
\end{alist}
